% GENERATED FILE. Edit canonical lessons, then run npm run generate:books.
% canonical-source-hash: fnv1a64:dfc8516de89f00c5
% canonical-lessons: MW-C59-seb, MW-C59-kelo, MW-C59-alu, MW-C59-tamatar, MW-C59-gajar

\chapter{Apple, Banana, Potato, Tomato, Carrot}
\label{ch:mw-apple-banana-potato-tomato-carrot}
\hlchaptermodality{59}

\begin{chapteropening}
\textbf{By the end of this chapter:} \emph{I can say an apple, a banana, a potato, a tomato and a carrot.}

Complete the last lesson of chapter 59: I can say an apple, a banana, a potato, a tomato and a carrot.
\end{chapteropening}

\section[\texorpdfstring{\mw{सेब} (seb)}{सेब (seb)}]{\mw{सेब} (seb) — an apple}

\label{lesson:MW-C59-seb}



\begin{quote}
Before the new one: say the Marwadi for a fish, then the Marwadi for rice.
\end{quote}

\subsection*{You'll want to know: \mw{सेब}}
\textbf{\mw{सेब}} — \emph{seb} — \textquotedblleft{}an apple\textquotedblright{}.

\begin{cousinweb}
\textbf{\mw{सेब}} (\emph{seb}) is a loan, usually traced to Persian \emph{sēb}, \textquotedblleft{}apple\textquotedblright{}.
\end{cousinweb}

\begin{grammarlens}[title={What you've built}]
One more word for this chapter.
\end{grammarlens}

\subsection*{Your turn}
\begin{itemize}
\raggedright
  \item \emph{Say it:} \emph{seb}
  \item \emph{Say it:} \emph{seb}, once more
  \item \emph{Say it:} say \emph{seb}
  \item \emph{Recall:} say the Marwadi for a fish, then the Marwadi for rice, then say \emph{seb} again
\end{itemize}

\begin{tcolorbox}[breakable,colback=teal!4,colframe=teal!35!black,title={Before you move on}]
What does \mw{सेब} mean? (\textquotedblleft{}An apple\textquotedblright{}.) Say it once more.
\end{tcolorbox}
\section[\texorpdfstring{\mw{केळो} (keḷo)}{केळो (keḷo)}]{\mw{केळो} (keḷo) — a banana}

\label{lesson:MW-C59-kelo}



\begin{quote}
Before the new one: say the Marwadi for rice, then the Marwadi for an apple.
\end{quote}

\subsection*{You'll want to know: \mw{केळो}}
\textbf{\mw{केळो}} — \emph{keḷo} — \textquotedblleft{}a banana\textquotedblright{}.

\begin{cousinweb}
\textbf{\mw{केळो}} (\emph{keḷo}) is a tadbhava, a word inherited from Sanskrit and reshaped by sound change. It comes from \emph{kadala}, \textquotedblleft{}the banana plant\textquotedblright{}. CDIAL gives no Marwari form in this entry (CDIAL 2712); the match is through Hindi \emph{kelā} and Gujarati \emph{keḷũ}. CDIAL takes the Sanskrit word itself from an Austroasiatic language.
\end{cousinweb}

\begin{grammarlens}[title={What you've built}]
One more word for this chapter.
\end{grammarlens}

\subsection*{Your turn}
\begin{itemize}
\raggedright
  \item \emph{Say it:} \emph{keḷo}
  \item \emph{Say it:} \emph{keḷo}, once more
  \item \emph{Say it:} say \emph{keḷo}
  \item \emph{Recall:} say the Marwadi for rice, then the Marwadi for an apple, then say \emph{keḷo} again
\end{itemize}

\begin{tcolorbox}[breakable,colback=teal!4,colframe=teal!35!black,title={Before you move on}]
What does \mw{केळो} mean? (\textquotedblleft{}A banana\textquotedblright{}.) Say it once more.
\end{tcolorbox}
\section[\texorpdfstring{\mw{आलू} (ālū)}{आलू (ālū)}]{\mw{आलू} (ālū) — a potato}

\label{lesson:MW-C59-alu}



\begin{quote}
Before the new one: say the Marwadi for an apple, then the Marwadi for a banana.
\end{quote}

\subsection*{You'll want to know: \mw{आलू}}
\textbf{\mw{आलू}} — \emph{ālū} — \textquotedblleft{}a potato\textquotedblright{}.

\begin{cousinweb}
\textbf{\mw{आलू}} (\emph{ālū}) is a tadbhava, a word inherited from Sanskrit and reshaped by sound change. It comes from \emph{ālu}, \textquotedblleft{}esculent root\textquotedblright{}. CDIAL gives no Marwari form in this entry (CDIAL 1388); the match is through Hindi \emph{ālū}. The potato took the name of an older edible root.
\end{cousinweb}

\begin{grammarlens}[title={What you've built}]
One more word for this chapter.
\end{grammarlens}

\subsection*{Your turn}
\begin{itemize}
\raggedright
  \item \emph{Say it:} \emph{ālū}
  \item \emph{Say it:} \emph{ālū}, once more
  \item \emph{Say it:} say \emph{ālū}
  \item \emph{Recall:} say the Marwadi for an apple, then the Marwadi for a banana, then say \emph{ālū} again
\end{itemize}

\begin{tcolorbox}[breakable,colback=teal!4,colframe=teal!35!black,title={Before you move on}]
What does \mw{आलू} mean? (\textquotedblleft{}A potato\textquotedblright{}.) Say it once more.
\end{tcolorbox}
\section[\texorpdfstring{\mw{टमाटर} (ṭamāṭar)}{टमाटर (ṭamāṭar)}]{\mw{टमाटर} (ṭamāṭar) — a tomato}

\label{lesson:MW-C59-tamatar}



\begin{quote}
Before the new one: say the Marwadi for a banana, then the Marwadi for a potato.
\end{quote}

\subsection*{You'll want to know: \mw{टमाटर}}
\textbf{\mw{टमाटर}} — \emph{ṭamāṭar} — \textquotedblleft{}a tomato\textquotedblright{}.

\begin{cousinweb}
\textbf{\mw{टमाटर}} (\emph{ṭamāṭar}) is a loan from English \emph{tomato}.
\end{cousinweb}

\begin{grammarlens}[title={What you've built}]
One more word for this chapter.
\end{grammarlens}

\subsection*{Your turn}
\begin{itemize}
\raggedright
  \item \emph{Say it:} \emph{ṭamāṭar}
  \item \emph{Say it:} \emph{ṭamāṭar}, once more
  \item \emph{Say it:} say \emph{ṭamāṭar}
  \item \emph{Recall:} say the Marwadi for a banana, then the Marwadi for a potato, then say \emph{ṭamāṭar} again
\end{itemize}

\begin{tcolorbox}[breakable,colback=teal!4,colframe=teal!35!black,title={Before you move on}]
What does \mw{टमाटर} mean? (\textquotedblleft{}A tomato\textquotedblright{}.) Say it once more.
\end{tcolorbox}
\section[\texorpdfstring{\mw{गाजर} (gājar)}{गाजर (gājar)}]{\mw{गाजर} (gājar) — a carrot}

\label{lesson:MW-C59-gajar}



\begin{quote}
Before the new one: say the Marwadi for a potato, then the Marwadi for a tomato.
\end{quote}

\subsection*{You'll want to know: \mw{गाजर}}
\textbf{\mw{गाजर}} — \emph{gājar} — \textquotedblleft{}a carrot\textquotedblright{}.

\begin{cousinweb}
\textbf{\mw{गाजर}} (\emph{gājar}) is a tadbhava, a word inherited from Sanskrit and reshaped by sound change. It comes from \emph{gārjara}, \textquotedblleft{}carrot\textquotedblright{}. CDIAL gives no Marwari form in this entry (CDIAL 4140); the match is through Hindi \emph{gājar} and Gujarati \emph{gājar}.
\end{cousinweb}

\begin{grammarlens}[title={What you've built}]
One more word for this chapter.
\end{grammarlens}

\subsection*{Your turn}
\begin{itemize}
\raggedright
  \item \emph{Say it:} \emph{gājar}
  \item \emph{Say it:} \emph{gājar}, once more
  \item \emph{Say it:} say \emph{gājar}
  \item \emph{Recall:} say the Marwadi for a potato, then the Marwadi for a tomato, then say \emph{gājar} again
\end{itemize}

\begin{tcolorbox}[breakable,colback=teal!4,colframe=teal!35!black,title={Before you move on}]
What does \mw{गाजर} mean? (\textquotedblleft{}A carrot\textquotedblright{}.) Say it once more.
\end{tcolorbox}
